\documentclass[10pt,a4paper]{amsart}
\usepackage{amsmath,amssymb,amsthm,graphicx}
% For commutative diagrams we use Xy-pic
\usepackage[all]{xy}

% We use 2cell for 2-commutative diagrams.
\xyoption{2cell}
% Theorem environments.
%
\theoremstyle{plain}
\newtheorem{theorem}[subsection]{Theorem}
\newtheorem{proposition}[subsection]{Proposition}
\newtheorem{lemma}[subsection]{Lemma}

\theoremstyle{definition}
\newtheorem{definition}[subsection]{Definition}
\newtheorem{example}[subsection]{Example}
\newtheorem{exercise}[subsection]{Exercise}
\newtheorem{situation}[subsection]{Situation}

\theoremstyle{remark}
\newtheorem{remark}[subsection]{Remark}
\newtheorem{remarks}[subsection]{Remarks}

\numberwithin{equation}{subsection}

% Macros
%
\def\lim{\mathop{\mathrm{lim}}\nolimits}
\def\colim{\mathop{\mathrm{colim}}\nolimits}
\def\Spec{\mathop{\mathrm{Spec}}}
\def\Hom{\mathop{\mathrm{Hom}}\nolimits}
\def\Ext{\mathop{\mathrm{Ext}}\nolimits}
\def\SheafHom{\mathop{\mathcal{H}\!\mathit{om}}\nolimits}
\def\SheafExt{\mathop{\mathcal{E}\!\mathit{xt}}\nolimits}
\def\Sch{\mathit{Sch}}
\def\Mor{\mathop{\mathrm{Mor}}\nolimits}
\def\Ob{\mathop{\mathrm{Ob}}\nolimits}
\def\Sh{\mathop{\mathit{Sh}}\nolimits}
\def\NL{\mathop{N\!L}\nolimits}
\def\CH{\mathop{\mathrm{CH}}\nolimits}
\def\proetale{{pro\text{-}\acute{e}tale}}
\def\etale{{\acute{e}tale}}
\def\QCoh{\mathit{QCoh}}
\def\Ker{\mathop{\mathrm{Ker}}}
\def\Im{\mathop{\mathrm{Im}}}
\def\Coker{\mathop{\mathrm{Coker}}}
\def\Coim{\mathop{\mathrm{Coim}}}

% Boxtimes
%
\DeclareMathSymbol{\boxtimes}{\mathbin}{AMSa}{"02}

%
% Macros for moduli stacks/spaces
%
\def\QCohstack{\mathcal{QC}\!\mathit{oh}}
\def\Cohstack{\mathcal{C}\!\mathit{oh}}
\def\Spacesstack{\mathcal{S}\!\mathit{paces}}
\def\Quotfunctor{\mathrm{Quot}}
\def\Hilbfunctor{\mathrm{Hilb}}
\def\Curvesstack{\mathcal{C}\!\mathit{urves}}
\def\Polarizedstack{\mathcal{P}\!\mathit{olarized}}
\def\Complexesstack{\mathcal{C}\!\mathit{omplexes}}
% \Pic is the operator that assigns to X its picard group, usage \Pic(X)
% \Picardstack_{X/B} denotes the Picard stack of X over B
% \Picardfunctor_{X/B} denotes the Picard functor of X over B
\def\Pic{\mathop{\mathrm{Pic}}\nolimits}
\def\Picardstack{\mathcal{P}\!\mathit{ic}}
\def\Picardfunctor{\mathrm{Pic}}
\def\Deformationcategory{\mathcal{D}\!\mathit{ef}}
\newenvironment{reference}{}{}
\newenvironment{slogan}{}{}
\newenvironment{history}{}{}
\usepackage{cleveref}
\crefname{originalproposition}{original proposition}{original propositions}
\Crefname{originalproposition}{Original proposition}{Original propositions}
\makeatletter
\newlabel{cohomology-proposition-vanishing-Noetherian}{{02UZ}{source}{}{sourcepointer-1}{}}
\newlabel{cohomology-proposition-vanishing-Noetherian@cref}{{[originalproposition][0][]02UZ}{[1][0][]source}{}{}{}}
\makeatother
\crefname{originallemma}{original lemma}{original lemmas}
\Crefname{originallemma}{Original lemma}{Original lemmas}
\makeatletter
\newlabel{lemma-regular-model-field}{{0C68}{source}{}{sourcepointer-2}{}}
\newlabel{lemma-regular-model-field@cref}{{[originallemma][0][]0C68}{[1][0][]source}{}{}{}}
\makeatother
\crefname{originallemma}{original lemma}{original lemmas}
\Crefname{originallemma}{Original lemma}{Original lemmas}
\makeatletter
\newlabel{lemma-numerical-type-minimal-model}{{0CA6}{source}{}{sourcepointer-3}{}}
\newlabel{lemma-numerical-type-minimal-model@cref}{{[originallemma][0][]0CA6}{[1][0][]source}{}{}{}}
\makeatother
\crefname{originallemma}{original lemma}{original lemmas}
\Crefname{originallemma}{Original lemma}{Original lemmas}
\makeatletter
\newlabel{lemma-divisor-special-fiber}{{0C5Z}{source}{}{sourcepointer-4}{}}
\newlabel{lemma-divisor-special-fiber@cref}{{[originallemma][0][]0C5Z}{[1][0][]source}{}{}{}}
\makeatother
\title{For a minimal model with coprime multiplicities the first cohomology of the special fibre surjects onto that of its reduction, with nontrivial kernel when they differ}
\author{The Stacks Project Authors}
\date{stacks-project:models.tex:local-snapshot-20261012; GNU FDL 1.2 or later}
\begin{document}
\maketitle
\noindent\textit{Editorial note.} The counted unit is the Stacks Project Lemma lemma-genus-reduction-smaller at source lines 5374-5383 together with its proof at source lines 5385-5429, a straight adjacent pairing in the chapter file. The chapter is the upstream file models.tex; the pool directory carries the longer name stacks-models and is not used as the chapter name anywhere in this card.\par\medskip
\noindent\textit{Editorial note.} One context span carries the regular-model situation at lines 4628-4637, which the retained statement names; it is context rather than a counted unit, and it does not overlap the counted interval.\par\medskip
\noindent\textit{Editorial note.} The retained proof quotes one proposition of the Cohomology of Schemes chapter and names two lemmas of this chapter, the one on the field of constants and the chain of subschemes, and the one comparing minimality of the model with minimality of its numerical type. All three are declared as non-load-bearing navigation pointers; the third is counted by a separate card of this lane and its interval is not reproduced here.\par\medskip
\noindent\textit{Editorial note.} The pointer to the vanishing proposition is the one external quotation of this card; its number 02UZ and the index line recorded for it were read off the shared tag index before assembly.\par\medskip
\noindent\textit{Editorial note.} No citation occurs in the retained spans, so no bibliography entry is supplied.\par\medskip
\noindent\textit{Editorial note.}  The retained situation carries a reference of its own, to the lemma writing the special fibre as a sum of its components with multiplicities, so that lemma is declared as a pointer as well: a context span that names a result must have its reference resolvable. The first revision omitted it and the engine reported an undefined reference.\par\medskip
\noindent\textit{Allowed standard prerequisites (editorial).}\begin{itemize}
\item The regular-model situation, minimality with coprime multiplicities, and vanishing of higher cohomology on the special fibre. The situation fixes a discrete valuation ring, a smooth projective curve over its fraction field, a regular proper model and the irreducible components of the special fibre with their multiplicities; minimality of the model is equivalent to minimality of the numerical type, and the hypothesis that the multiplicities have greatest common divisor one is what makes the reduction of the special fibre differ from the fibre only when some multiplicity exceeds one. Cohomology of coherent sheaves on the special fibre vanishes in degrees at least two, so a surjection of sheaves induces a surjection on first cohomology. The author's own lemmas of the chapter and the quoted vanishing proposition are the allowed inputs.. The counted lemma is one half of the genus formula of the section: it bounds the difference between the genus of the special fibre and that of its reduction from one side, the companion lemma giving the other side.
\end{itemize}\medskip
\ifdefined\hypertarget\hypertarget{sourcepointer-1}{}\fi
\noindent\textit{Source reference.} Original proposition 02UZ: focused-seeds/AGP/sources/stacks-snapshot-20261011/tags.txt, source lines 3750--3750. Allowed input of the Cohomology of Schemes chapter quoted by the retained proof for the vanishing of cohomology in degrees at least two on the special fibre; the pointer is registered so the reference resolves and the cited statement is not reproduced inside the span. The number is the stable Stacks tag of the target, 02UZ, looked up in the shared tag index.\par\medskip
\ifdefined\hypertarget\hypertarget{sourcepointer-2}{}\fi
\noindent\textit{Source reference.} Original lemma 0C68: focused-seeds/AGP/sources/stacks-snapshot-20261011/tags.txt, source lines 15288--15288. Navigation pointer invoked by the retained proof for the chain from the reduction of the special fibre to the fibre itself and for the parallel short exact sequence; that lemma is an allowed input of the same section and its proof is not reproduced here. The number is the stable Stacks tag of the target, looked up in the shared tag index.\par\medskip
\ifdefined\hypertarget\hypertarget{sourcepointer-3}{}\fi
\noindent\textit{Source reference.} Original lemma 0CA6: focused-seeds/AGP/sources/stacks-snapshot-20261011/tags.txt, source lines 15426--15426. Navigation pointer invoked by the retained proof in its use of minimality; that lemma is counted by a separate card of this lane, so its proof is not reproduced here. The number is the stable Stacks tag of the target, looked up in the shared tag index.\par\medskip
\ifdefined\hypertarget\hypertarget{sourcepointer-4}{}\fi
\noindent\textit{Source reference.} Original lemma 0C5Z: focused-seeds/AGP/sources/stacks-snapshot-20261011/tags.txt, source lines 15279--15279. Navigation pointer carried by the retained situation, which names this lemma when it writes the special fibre as a sum of its components with multiplicities; the situation is context here, so its own reference must resolve. The number is the stable Stacks tag of the target, looked up in the shared tag index.\par\medskip
\begin{situation}
\label{situation-regular-model}
Let $R$ be a discrete valuation ring with fraction field $K$,
residue field $k$, and uniformizer $\pi$.
Let $C$ be a smooth projective curve over $K$ with $H^0(C, \mathcal{O}_C) = K$.
Let $X$ be a regular proper model of $C$.
Let $C_1, \ldots, C_n$ be the irreducible components of the special
fibre $X_k$. Write $X_k = \sum m_i C_i$ as in
Lemma \ref{lemma-divisor-special-fiber}.
\end{situation}

\begin{lemma}
\label{lemma-genus-reduction-smaller}
In Situation \ref{situation-regular-model} assume $X$ is a minimal model,
$\gcd(m_1, \ldots, m_n) = 1$, and $H^0((X_k)_{red}, \mathcal{O}) = k$. Then
the map
$$
H^1(X_k, \mathcal{O}_{X_k}) \to H^1((X_k)_{red}, \mathcal{O}_{(X_k)_{red}})
$$
is surjective and has a nontrivial kernel as soon as $(X_k)_{red} \not = X_k$.
\end{lemma}

\begin{proof}
By vanishing of cohomology in degrees $\geq 2$ over $X_k$
(Cohomology, Proposition \ref{cohomology-proposition-vanishing-Noetherian})
any surjection of abelian sheaves on $X_k$ induces a surjection on $H^1$.
Consider the sequence
$$
(X_k)_{red} = Z_0 \subset Z_1 \subset \ldots \subset Z_m = X_k
$$
of Lemma \ref{lemma-regular-model-field}. Since the field maps
$H^0(Z_j, \mathcal{O}_{Z_j}) \to
H^0((X_k)_{red}, \mathcal{O}_{(X_k)_{red}}) = k$
are injective we conclude that $H^0(Z_j, \mathcal{O}_{Z_j}) = k$ for
$j = 0, \ldots, m$. It follows that
$H^0(X_k, \mathcal{O}_{X_k}) \to H^0(Z_{m - 1}, \mathcal{O}_{Z_{m - 1}})$
is surjective. Let $C = C_{i_m}$. Then $X_k = Z_{m - 1} + C$.
Let $\mathcal{L} = \mathcal{O}_X(-Z_{m - 1})|_C$.
Then $\mathcal{L}$ is an invertible $\mathcal{O}_C$-module.
As in the proof of Lemma \ref{lemma-regular-model-field}
there is an exact sequence
$$
0 \to \mathcal{L} \to \mathcal{O}_{X_k} \to \mathcal{O}_{Z_{m - 1}} \to 0
$$
of coherent sheaves on $X_k$. We conclude that
we get a short exact sequence
$$
0 \to
H^1(C, \mathcal{L}) \to H^1(X_k, \mathcal{O}_{X_k}) \to
H^1(Z_{m - 1}, \mathcal{O}_{Z_{m - 1}}) \to 0
$$
The degree of $\mathcal{L}$ on $C$ over $k$ is
$$
(C \cdot -Z_{m - 1}) = (C \cdot C - X_k) = (C \cdot C)
$$
Set $\kappa = H^0(C, \mathcal{O}_C)$ and $w = [\kappa : k]$.
By definition of the degree of an invertible sheaf we see that
$$
\chi(C, \mathcal{L}) =
\chi(C, \mathcal{O}_C) + (C \cdot C) =
w(1 - g_C) + (C \cdot C)
$$
where $g_C$ is the genus of $C$. This expression is $< 0$ as $X$ is minimal
and hence $C$ is not an exceptional curve of the first kind
(see proof of Lemma \ref{lemma-numerical-type-minimal-model}).
Thus $\dim_k H^1(C, \mathcal{L}) > 0$ which finishes the proof.
\end{proof}

\end{document}
